%%%%
% MTecknology's Resume
%%%%
% Author: Michael Lustfield
% License: CC-BY-4
% - https://creativecommons.org/licenses/by/4.0/legalcode.txt
%%%%

\documentclass[letterpaper,10pt]{article}
%%%%%%%%%%%%%%%%%%%%%%%
%% BEGIN_FILE: mteck.sty
%% NOTE: Everything between here and END_FILE can
%% be relocated to "mteck.sty" and then included with:
%\usepackage{mteck}

% Dependencies
% NOTE: Some packages (lastpage, hyperref) require second build!
\usepackage[empty]{fullpage}
\usepackage{titlesec}
\usepackage{enumitem}
\usepackage[hidelinks]{hyperref}
\usepackage{fancyhdr}
\usepackage{fontawesome5}
\usepackage{multicol}
\usepackage{bookmark}
\usepackage{lastpage}

% Sans-Serif
%\usepackage[sfdefault]{FiraSans}
%\usepackage[sfdefault]{roboto}
%\usepackage[sfdefault]{noto-sans}
%\usepackage[default]{sourcesanspro}

% Serif
\usepackage{CormorantGaramond}
\usepackage{charter}

% Colors
% Use with \color{Name}
% Can wrap entire heading with {\color{accent} [...] }
\usepackage{xcolor}
\definecolor{accentTitle}{HTML}{003366}
\definecolor{accentText}{HTML}{003366}
\definecolor{accentLine}{HTML}{003366}

% Misc. Options
\pagestyle{fancy}
\fancyhf{}
\fancyfoot{}
\renewcommand{\headrulewidth}{0pt}
\renewcommand{\footrulewidth}{0pt}
\urlstyle{same}

% Adjust Margins
\addtolength{\oddsidemargin}{-0.7in}
\addtolength{\evensidemargin}{-0.5in}
\addtolength{\textwidth}{1.19in}
\addtolength{\topmargin}{-0.7in}
\addtolength{\textheight}{1.4in}

\setlength{\multicolsep}{-3.0pt}
\setlength{\columnsep}{-1pt}
\setlength{\tabcolsep}{0pt}
\setlength{\footskip}{3.7pt}
\raggedbottom
\raggedright

% ATS Readability
\input{glyphtounicode}
\pdfgentounicode=1

%-----------------%
% Custom Commands %
%-----------------%

% Centered title along with subtitle between HR
% Usage: \documentTitle{Name}{Details}
\newcommand{\documentTitle}[2]{
  \begin{center}
    {\Huge\color{accentTitle} #1}
    \vspace{10pt}
    {\color{accentLine} \hrule}
    \vspace{2pt}
    %{\footnotesize\color{accentTitle} #2}
    \footnotesize{#2}
    \vspace{2pt}
    {\color{accentLine} \hrule}
  \end{center}
}

% Create a footer with correct padding
% Usage: \documentFooter{\thepage of X}
\newcommand{\documentFooter}[1]{
  \setlength{\footskip}{10.25pt}
  \fancyfoot[C]{\footnotesize #1}
}

% Simple wrapper to set up page numbering
% Usage: \numberedPages
% WARNING: Must run pdflatex twice!
\newcommand{\numberedPages}{
  \documentFooter{\thepage/\pageref{LastPage}}
}

% Section heading with horizontal rule
% Usage: \section{Title}
\titleformat{\section}{
  \vspace{-5pt}
  \color{accentText}
  \raggedright\large\bfseries
}{}{0em}{}[\color{accentLine}\titlerule]

% Section heading with separator and content on same line
% Usage: \tinysection{Title}
\newcommand{\tinysection}[1]{
  \phantomsection
  \addcontentsline{toc}{section}{#1}
  {\large{\bfseries\color{accentText}#1} {\color{accentLine} |}}
}

% Indented line with arguments left/right aligned in document
% Usage: \heading{Left}{Right}
\newcommand{\heading}[2]{
  \hspace{10pt}#1\hfill#2\\
}

% Adds \textbf to \heading
\newcommand{\headingBf}[2]{
  \heading{\textbf{#1}}{\textbf{#2}}
}

% Adds \textit to \heading
\newcommand{\headingIt}[2]{
  \heading{\textit{#1}}{\textit{#2}}
}

% Template for itemized lists
% Usage: \begin{resume_list} [items] \end{resume_list}
\newenvironment{resume_list}{
  \vspace{-7pt}
  \begin{itemize}[itemsep=-2px, parsep=1pt, leftmargin=30pt]
}{
  \end{itemize}
  %\vspace{-2pt}
}

% Formats an item to use as an itemized title
% Usage: \itemTitle{Title}
\newcommand{\itemTitle}[1]{
  \item[] \underline{#1}\vspace{4pt}
}

% Bullets used in itemized lists
\renewcommand\labelitemi{--}

%% END_FILE: mteck.sty
%%%%%%%%%%%%%%%%%%%%%%


%===================%
% John Doe's Resume %
%===================%

%\numberedPages % NOTE: lastpage requires a second build
%\documentFooter{\thepage of 2} % Does similar without using lastpage
\begin{document}

  %---------%
  % Heading %
  %---------%

%---------%
  % Heading %
  %---------%

%---------%
  % Heading %
  %---------%

  \documentTitle{LA NHAT HY \\ \vspace{0.5em} \large Embedded Linux Software Engineer}{
    \href{tel:0836476472}{
      \raisebox{-0.05\height} \faPhone\ (+84) 836 476 472} ~ | ~
    \href{mailto:lahyus7399@gmail.com}{
      \raisebox{-0.15\height} \faEnvelope\ lahyus7399@gmail.com} ~ | ~
    \href{https://github.com/LahyUS}{
      \raisebox{-0.15\height} \faGithub\ github.com/LahyUS}  ~ | ~
    \href{https://www.linkedin.com/in/nhat-hy-la-40b379221/}{
      \raisebox{-0.15\height} \faLinkedin\ linkedin.com/in/nhat-hy-la-40b379221/}
  }

  %---------%
  % Summary %
  %---------%

  \tinysection{Summary}
  \textbf{I bridge the gap between Hardware and Intelligence.} As a Software Engineer specializing in \textbf{Software Development}, \textbf{Embedded Systems}, and \textbf{Applied AI}, I build systems that are both robust and smart. My expertise spans from low-level C/C++ and Linux optimizations to deploying Computer Vision and RAG solutions on scalable platforms. I am committed to evidence-based engineering, focusing on precision, clarity, and documentable results.

  %--------%
  % Skills %
  %--------%

  \section{\faTools\ SKILLS}
 
   \begin{itemize}[itemsep=1pt, parsep=1pt, leftmargin=15pt, label={}]
     \item \textbf{\faCode\ Languages}: C11, C++17, Python 3.x, Bash, SQL, JavaScript (HTMX).
     \item \textbf{\faMicrochip\ Embedded}: ARM Linux (Realtek/Spider), Buildroot, Xen Hypervisor, CMake, POSIX Threads, NEON SIMD.
     \item \textbf{\faGlobe\ Protocols}: WebRTC, ONVIF (Profile S/T), MQTT, RTSP, SSL/TLS (OpenSSL/DTLS), WebSocket, mDNS (Avahi).
     \item \textbf{\faBrain\ AI \& Data}: RAG (LangChain), Computer Vision (OpenCV), Deep Learning (PyTorch), PaddleOCR.
     \item \textbf{\faLayerGroup\ Frameworks \& Libs}: FFmpeg, SDL2, libopus, libdatachannel, Flask, MySQL, SQLite, Matplotlib.
     \item \textbf{\faToolbox\ Tools}: Git (GitHub/GitLab), Jira, Confluence, GDB, Valgrind, Docker, VS Code, Tkinter, Pytest.
   \end{itemize}
 
   %-----------%
   % Education %
   %-----------%
 
   \section{\faGraduationCap\ EDUCATION}
 
   \headingBf{VNU-HCM University of Science}{Aug 2018 -- Sep 2024}
   \headingIt{Bachelor of Computer Science}{GPA: 7.7/10}
   \headingIt{Major: Computer Vision}{}
 
   \section{\faCertificate\ CERTIFICATIONS}
   \begin{resume_list}
     \item TOEIC Listening - Reading: 825 pts | Speaking - Writing: 290 pts
     \item Machine Learning Specialization (Coursera/DeepLearning.AI) -- \href{https://coursera.org/verify/specialization/KLJQLHVZTC92}{\underline{View Credential}}
   \end{resume_list}
 
   %------------%
   % Experience %
   %------------%
 
%------------%
  % Experience %
  %------------%

%------------%
  % Experience %
  %------------%

  \section{\faBriefcase\ WORK EXPERIENCE}
   
   \headingBf{Ampacs Corporation}{04/2025 -- Present}
   \headingIt{Embedded Software Engineer (R\&D)}{}
   \begin{resume_list}
     \item \textbf{\faSitemap\ Project: AmpacsSDK (Modular IoT Framework)}
       \begin{itemize}[itemsep=1pt, leftmargin=15pt, label=-]
         \item \textbf{Description}: A modular IoT Framework SDK designed for embedded devices, providing backend-agnostic MQTT connectivity and simplified HAL abstractions.
         \item \textbf{Role}: Core SDK Developer
         \item \textbf{Key Contributions}:
           \begin{itemize}[itemsep=0pt, leftmargin=10pt, label=$\circ$]
              \item Decoupled MQTTS library from specific backends (Paho/EMQX) using opaque handles.
              \item Implemented a dynamic callback manager for runtime registration of commands.
           \end{itemize}
         \item \textbf{Achievement}: Enabled seamless vendor switching without code changes.
         \item \textbf{Tech Stack}: C, MQTT, Paho, PlantUML, Linux.
       \end{itemize}

      \item \textbf{\faMicrophone\ Project: Low-Latency VoIP \& Opus Codec Integration}
        \begin{itemize}[itemsep=1pt, leftmargin=15pt, label=-]
          \item \textbf{Description}: Integration of the high-performance Opus (libopus) audio codec into the IP Camera’s ARMv7 firmware, replacing legacy codecs with voice-optimized, low-latency VoIP capabilities.
          \item \textbf{Role}: Primary Audio Engineer
          \item \textbf{Key Contributions}:
            \begin{itemize}[itemsep=0pt, leftmargin=10pt, label=$\circ$]
               \item \textbf{CPU Optimization}: Achieved a 48\% reduction in CPU overhead (\textbf{35\% to 18\%}) by enabling fixed-point math and leveraging ARM NEON (SIMD) optimizations for the Cortex-A7 architecture.
               \item \textbf{Pipeline Engineering}: Restored the complete \textbf{Resampler $\rightarrow$ Mixer $\rightarrow$ Playback} pipeline, enabling multi-source mixing (voice + ringtones) and resolving legacy audio glitching.
               \item \textbf{Build Orchestration}: Migrated build integration to a Buildroot package workflow, utilizing \textbf{PkgConfig} for sysroot discovery and removing hardcoded toolchain dependencies.
               \item \textbf{Synchronization Logic}: Engineered a monolithic timestamping mechanism and buffer profile injection to ensure jitter-free, 20ms frame pacing within the Realtek SDK buffers.
            \end{itemize}
          \item \textbf{Achievement}: Delivered crystal-clear two-way audio with professional-grade lip-sync accuracy and a significantly reduced system resource footprint.
          \item \textbf{Tech Stack}: C, libopus, Buildroot, NEON SIMD, POSIX Threads, CMake, PkgConfig.
        \end{itemize}

      \item \textbf{\faLock\ Project: Secure Snapshot Encryption (AES-128-CBC)}
        \begin{itemize}[itemsep=1pt, leftmargin=15pt, label=-]
          \item \textbf{Description}: A hardware-bound encryption module for IP camera snapshots, ensuring user privacy by encrypting JPEG data using keys unique to the device’s hardware (UDID).
          \item \textbf{Role}: Primary Security Developer
          \item \textbf{Key Contributions}:
            \begin{itemize}[itemsep=0pt, leftmargin=10pt, label=$\circ$]
               \item \textbf{API Design}: Developed a C-based AES-128-CBC encryption module, exposing a \textbf{clean, high-level API} (\texttt{image\_encrypt\_file/buffer}) that decoupled complex AES-128-CBC/PKCS#7 logic from the primary application.
               \item \textbf{Hardware-Bound Security}: Implemented a deterministic key derivation system based on the device’s UDID, ensuring that cryptographic materials are never stored on disk and are zeroed in memory after use.
               \item \textbf{Multi-Mode Support}: Engineered both \textbf{Standard Mode} (randomly salted IVs) and \textbf{Legacy Mode} (unsalted), providing flexibility for different server-side decryption requirements.
               \item \textbf{Verification Suite}: Built a comprehensive C test application and an interactive \textbf{OpenSSL cross-check} tool to validate server-side decryption and ensure 100\% data integrity.
            \end{itemize}
          \item \textbf{Achievement}: Delivered an intuitive, production-grade security library used by the entire team, protecting snapshots in transit with zero measurable performance impact.
          \item \textbf{Tech Stack}: C, AES-CBC, libcurl, CMake, OpenSSL.
        \end{itemize}

       \item \textbf{\faMoon\ Project: HAL-Driven Autonomous Day/Night Controller}
        \begin{itemize}[itemsep=1pt, leftmargin=15pt, label=-]
          \item \textbf{Description}: A production-grade lighting management system for IP cameras, enabling autonomous transitions between high-fidelity color and IR-illuminated monochrome modes based on environmental light levels.
          \item \textbf{Role}: System Software Engineer
          \item \textbf{Key Contributions}:
            \begin{itemize}[itemsep=0pt, leftmargin=10pt, label=$\circ$]
               \item \textbf{HAL Architecture}: Designed and implemented a Hardware Abstraction Layer (HAL) using C function pointers to decouple core logic from the Realtek SDK, ensuring 100\% portability to new SoC vendors.
               \item \textbf{Autonomous State Machine}: Engineered a robust, thread-safe asynchronous state machine to monitor light sensors and synchronize binary triggers (ISP mode, IR-cut filter, and IR LED) with zero-flicker transitions.
               \item \textbf{Automated Test Harness}: Developed a \textbf{Mock HAL} for x86 environments, allowing for 100\% logic verification and regression testing on non-embedded systems, reducing hardware-dependency during development.
               \item \textbf{State Synchronization}: Integrated hardware transitions with ONVIF imaging standards and NVRAM, ensuring persistent state management and consistent behavior across network reboots.
            \end{itemize}
          \item \textbf{Achievement}: Delivered a highly portable, fail-safe system with synchronized hardware states, passing all ONVIF conformance tests for Profile S/T.
          \item \textbf{Tech Stack}: C, Linux, Realtek SDK, POSIX Threads, Buildroot, CMake, Shell Scripting.
        \end{itemize}

      \item \textbf{\faVideo\ Project: Autonomous WebRTC Streaming \& Discovery System}
        \begin{itemize}[itemsep=1pt, leftmargin=15pt, label=-]
          \item \textbf{Description}: Architected a standalone P2P streaming solution enabling real-time H.264 video delivery directly from the camera to web browsers without external server dependencies.
          \item \textbf{Role}: Streaming Software Engineer
          \item \textbf{Key Contributions}:
            \begin{itemize}[itemsep=0pt, leftmargin=10pt, label=$\circ$]
               \item \textbf{Embedded Signaling}: Developed an in-camera \textbf{WebSocket signaling server} using C++17 and libdatachannel, eliminating the need for cloud-based signaling in local network deployments.
               \item \textbf{Zero-Config Discovery}: Integrated \textbf{mDNS (Avahi)} with dynamic port resolution and multi-homed routing, enabling automatic camera discovery and path selection across complex subnets.
               \item \textbf{Performance Optimization}: Migrated mDNS listener loops to a \textbf{select()}-based architecture, eliminating 100\% CPU spin on embedded ARM hardware while maintaining high responsiveness.
               \item \textbf{Race Condition Resilience}: Engineered a thread-safe \textbf{ICE candidate queuing system} in C++ to handle signaling race conditions and implemented client-side connection racing for near-instant connectivity.
            \end{itemize}
          \item \textbf{Achievement}: Achieved near-instant video playback (under 1 second) by building a standalone streaming host that eliminates the need for expensive external signaling servers. Also serves as a flexible architecture supporting Local, Product, and Debug modes.
          \item \textbf{Tech Stack}: C++17, WebRTC, WebSocket, mDNS (Avahi), Python, H.264, CMake.
        \end{itemize}
   
     \item \textbf{\faNetworkWired\ Project: ONVIF Protocol Stack}
       \begin{itemize}[itemsep=1pt, leftmargin=15pt, label=-]
         \item \textbf{Description}: Implementation of ONVIF Profile S/T standards ensuring interoperability for IP Cameras.
         \item \textbf{Role}: Protocol Engineer
         \item \textbf{Key Contributions}:
           \begin{itemize}[itemsep=0pt, leftmargin=10pt, label=$\circ$]
              \item \textbf{Day/Night Integration}: Mapped ONVIF imaging commands to trigger the system's DNM engine, allowing Auto/Manual mode switching via network clients.
              \item Resolved critical profile lookup bugs in \texttt{GetImagingSettings}.
              \item Synchronized IR cut filter state between ONVIF, ISP, and NVRAM.
           \end{itemize}
         \item \textbf{Achievement}: Passed ONVIF conformance tests for Profile S/T.
         \item \textbf{Tech Stack}: C++, gSOAP, ONVIF, Python.
       \end{itemize}

 

 
      \item \textbf{\faPlayCircle\ Project: Secure Media Player (SSL-FFPlayer)}
       \begin{itemize}[itemsep=1pt, leftmargin=15pt, label=-]
         \item \textbf{Description}: A multi-threaded, securely encrypted media player designed for real-time streaming with strict A/V synchronization.
         \item \textbf{Role}: Multimedia Software Engineer
         \item \textbf{Key Contributions}:
           \begin{itemize}[itemsep=0pt, leftmargin=10pt, label=$\circ$]
              \item \textbf{Architecture}: Designed a decoupled 6-thread pipeline (Network $\rightarrow$ Demux $\rightarrow$ Decode $\rightarrow$ Render) using thread-safe queues and ring buffers to isolate network I/O from playback.
              \item \textbf{A/V Sync}: Engineered a master-slave synchronization logic where video frames are slaved to the audio clock (counted from hardware callbacks), eliminating drift.
              \item \textbf{Push-Pull Audio}: Implemented a lock-free ring buffer model to decouple the decoding thread ("Push") from the SDL audio callback ("Pull"), preventing stuttering.
              \item \textbf{Security}: Integrated OpenSSL/DTLS directly into the transport layer for end-to-end encrypted streaming.
           \end{itemize}
         \textbf{Achievement}: Built a production-grade secure player from scratch, handling jitter and latency autonomously.
         \item \textbf{Tech Stack}: C, FFmpeg, SDL2, OpenSSL, Pthreads.
       \end{itemize}
   \end{resume_list}

  \vspace{5pt}\hrule\vspace{5pt}

  \headingBf{Renesas Design Vietnam}{08/2022 -- 05/2024}
   \headingIt{Embedded Software Engineer}{}
   \begin{resume_list}
     \item \textbf{\faSync\ Project: OTA Installer (SOTA/FOTA)}
       \begin{itemize}[itemsep=1pt, leftmargin=15pt, label=-]
         \item \textbf{Description}: A fail-safe Firmware Over-The-Air installer for Spider S4/S4N automotive boards (Linux/Xen).
         \item \textbf{Role}: Embedded Software Engineer
         \item \textbf{Key Contributions}:
            \begin{itemize}[itemsep=0pt, leftmargin=10pt, label=$\circ$]
              \item Implemented robust fallback mechanisms to revert to stable firmware upon failure.
              \item Designed custom binary data frames (Checksum, Length) for reliability.
              \item Wrote Shell scripts to orchestrate updates for dual targets (eMMC/Flash).
            \end{itemize}
         \item \textbf{Achievement}: Delivered a fully automated, fail-safe update mechanism for automotive grade hardware.
         \item \textbf{Tech Stack}: C, Shell Script, Linux, Xen Hypervisor.
       \end{itemize}



 
      \item \textbf{\faRobot\ Project: AI Chatbot for MCAL (RAG System)}
        \begin{itemize}[itemsep=1pt, leftmargin=15pt, label=-]
          \item \textbf{Description}: An AI assistant for MCAL (Microcontroller Abstraction Layer) developers to query complex AUTOSAR specifications and C source code using domain-aware RAG.
          \item \textbf{Role}: AI Software Engineer
          \item \textbf{Key Contributions}:
             \begin{itemize}[itemsep=0pt, leftmargin=10pt, label=$\circ$]
               \item \textbf{Architecture}: Architected a decoupled microservices system separating Model, Document, and DB servers, enabling multi-node deployment via secure tunneling (ngrok).
               \item \textbf{RAG Pipeline}: Engineered a robust Retrieval-Augmented Generation pipeline using \textbf{LangChain} and \textbf{FAISS/ChromaDB}, leveraging quantized 13B models for local high-accuracy inference.
               \item \textbf{Fine-Tuning}: Leveraged Parameter-Efficient Fine-Tuning (\textbf{PEFT}) technique, specifically \textbf{LoRA}, to adapt base LLMs to specialized AUTOSAR embedded standards and technical C syntax.
               \item \textbf{Document Engineering}: Developed a structured data extraction engine to parse unstructured PDFs and C headers into optimized middle-data JSON for semantic indexing.
               \item \textbf{Optimization}: Integrated \textbf{Streaming Responses} using FastAPI and implemented a secure session heartbeat mechanism with automated cleanup logic.
             \end{itemize}
          \item \textbf{Achievement}: Significantly reduced specification lookup time for developers by delivering context-bounded accuracy (passing 100\% of internal technical test scenarios).
          \item \textbf{Tech Stack}: Python, LangChain, FAISS, LoRA, Transformers, FastAPI, Flask, Gradio, MySQL.
        \end{itemize}
   \end{resume_list}
 
   %----------------------------%
   % Projects %
    %----------------------------%
    \section{\faLaptopCode\ FREELANCE PROJECTS}
  
   \headingBf{ClinicPRO: Clinic Management Platform (v2 Web)}{11/2025 -- Present}
   \begin{resume_list}
     \item \textbf{Description}: A high-performance, local Single-Page Application (SPA) designed for dental clinics, integrating clinical documentation, real-time scheduling, and granular financial auditing.
     \item \textbf{Role}: Full Stack Developer (Solo)
     \item \textbf{Key Contributions}:
        \begin{itemize}[itemsep=0pt, leftmargin=10pt, label=$\circ$]
           \item \textbf{Architecture}: Architected a reactive **SPA experience** using **Flask and HTMX**, eliminating page reloads and providing a modern desktop-app feel with minimal client-side overhead.
           \item \textbf{Clinical Systems}: Engineered an **interactive 32-tooth dental chart** (JSON-backed) and a 4-status **Kanban board** synchronized with **FullCalendar** for end-to-end patient workflow management.
           \item \textbf{Financial Engine}: Developed a **partial payment tracking system** with automated debt calculation and clinic expense logging, supporting complex multi-session billing cycles.
           \item \textbf{Security \& Data}: Implemented a **15+ table normalized SQLite schema** with 3-level **RBAC** (Owner, Doctor, Receptionist) and integrated automated database backups for local data sovereignty.
           \item \textbf{Integration}: Leveraged **PaddleOCR** for automated patient record digitization and **Chart.js** to visualize clinic growth trends and treatment statistics.
        \end{itemize}
     \item \textbf{Achievement}: Delivered a portable, production-ready solution (bundled via **PyInstaller**) facilitating complete clinical and administrative operations for local dental practices.
     \item \textbf{Tech Stack}: Python (Flask), HTMX, Bootstrap 5, SQLite, Chart.js, PaddleOCR, PyInstaller, Pytest.
   \end{resume_list}
 
   \section{\faGraduationCap\ ACADEMIC PROJECTS}

   \headingBf{Video Anomaly Detection (Thesis)}{01/2022 -- 07/2022}
   \begin{resume_list}
     \item \textbf{Description}: Research and system integration of deep learning architectures for abnormal event detection in video surveillance data.
     \item \textbf{Role}: DL Researcher \& Integration Engineer
     \item \textbf{Key Contributions}:
        \begin{itemize}[itemsep=0pt, leftmargin=10pt, label=$\circ$]
           \item \textbf{SOTA Research}: Performed a comprehensive survey of state-of-the-art anomaly detection methodologies, focused on Transformer-based and GNN architectures for time-series modeling.
           \item \textbf{System Integration}: Architected the software integration layer to bridge complex deep learning approaches into functional surveillance frameworks, ensuring robust runtime data flow.
           \item \textbf{Model Evaluation}: Conducted rigorous training and evaluation cycles on benchmark datasets (UCSD Ped2, CUHK Avenue), optimizing hyperparameters and binary detection thresholds.
           \item \textbf{Visualization}: Engineered a pixel-wise \textbf{Image Similarity Module} to quantify and visualize anomalies through reconstruction error heatmaps and difference-mask generation.
        \end{itemize}
     \item \textbf{Achievement}: Successfully integrated advanced DL research into a functional surveillance architecture, providing meaningful visual intuition for abnormal event localization.
     \item \textbf{Tech Stack}: Python, PyTorch (CUDA), OpenCV, Matplotlib, Numpy.
   \end{resume_list}


\end{document}